% Lösungen Einheit 4 von 5 – Parameter aus Bedingungen bestimmen (zusammenbau v0.1)
\einheitenkopf{Einheit 4 von 5 · Parameter aus Bedingungen bestimmen}

\erg{16}{a) vorwärts \quad b) Nullstellen $0$ und $a$; $A = \left[-\frac{x^3}{3} + \frac{a x^2}{2}\right]_0^a$, also $A = \frac{1}{6}a^3$. \quad c) Integral $H_k(2) - H_k(0) = 0 - 4k = -4k$; das Flächenstück liegt unter der x-Achse, also $4k = 1$ und $k = 0{,}25$. \quad d) 20 cm sind 4 LE, also die Grenzen $\pm 2$; 1 FE entspricht 25 cm², also sind 400 cm² genau 16 FE; $c \cdot \frac{32}{3} = 16$, also $c = 1{,}5$. \quad e) $f_k'(x) = 4x(x^2 - k)$, Tiefstellen $\pm\sqrt{k}$; $f_k(\sqrt{k}) = -k^2 = -9$ mit $k > 0$ liefert $k = 3$. \quad f) $k > 5$: Quadrieren ($g(x) > 0$) liefert $k x^2 + 100 = \left(\frac{x^2}{4} + 10\right)^2$, also $x^2\left(k - 5 - \frac{x^2}{16}\right) = 0$; neben $x = 0$ gibt es genau für $k > 5$ die Lösungen $\pm 4\sqrt{k - 5}$. \quad g) Sprungfrei ist er schon: $f(0) = 2 = p_a(0)$. Knickfrei: $f'(0) = 2e^0 = 2$ und $p_a'(0) = b$, also $b = 2$. \quad h) Boden: $h_k(r_0) = 0$, $r_0^2 = \frac{8}{k}$; Rand bei $y = 3$ (30 cm): $r_1^2 = \frac{8}{k - 3}$; aus $r_1 = 2r_0$ folgt $\frac{8}{k - 3} = \frac{32}{k}$, also $k = 4$. \quad i) $f_a'(x) = -2ax + 2 = 0$: $v = \frac{1}{a}$, $f_a(v) = \frac{1}{a} = v$. Trapez mit den parallelen Seiten $2v$ und $v$ und der Höhe $v$: $\frac{2v + v}{2} \cdot v = 1{,}5v^2 = 24$, also $v = 4$ und $a = 0{,}25$.}
\erg{17}{a) Das Flächenstück liegt unter der x-Achse, das Integral $-\frac{a^3}{6}$ ist negativ; der Inhalt ist $\frac{a^3}{6} = \frac{4}{3}$, also $a^3 = 8$ und $a = 2$. \quad b) Der Flächeninhalt lässt sich als Term in $a$ berechnen; ihn gleich $10$ zu setzen ergibt eine Gleichung, deren einzige Unbekannte $a$ ist. \quad c) Nullstellen $0$ und $6$; $A = a \cdot \left[3x^2 - \frac{x^3}{3}\right]_0^6 = 36a = 18$, also $a = 0{,}5$; größte Ausdehnung $b_{0{,}5}(3) = 4{,}5$ m.}
